\section{GEMM vs. GEMV}

\begin{table}[H]
\centering
\caption{Comparación de las características de cómputo entre Prefill y Decode}
\begin{tabular}{lcc}
\toprule
& \textbf{Prefill} & \textbf{Decode} \\
\midrule
Operación central & GEMM (matriz $\times$ matriz) & GEMV (vector $\times$ matriz) \\
Longitud de Q & $N$ (longitud del prompt) & 1 \\
Tipo de cuello de botella & Intensivo en cómputo & Intensivo en ancho de banda de memoria \\
Uso de la GPU & Aprovechamiento pleno del cómputo & Cómputo ocioso en gran medida \\
Dirección de optimización & Más FLOPs & Mayor ancho de banda de memoria \\
Característica de latencia & TTFT (latencia del primer token) & TPS (tokens por segundo) \\
\bottomrule
\end{tabular}
\end{table}

\begin{warningbox}{El dilema de eficiencia de la fase de Decode}
En la fase de Decode se genera solo un token por paso, pero es necesario cargar todo el peso del modelo y el KV-Cache. El cómputo de la GPU queda ocioso en gran medida y el ancho de banda de memoria se convierte en el cuello de botella. Por eso la arquitectura de memoria unificada del Apple M5 Max (de alto ancho de banda) destaca en la fase de Decode.
\end{warningbox}
